\section{Семинар \thesection}
\subsection{Задачи по теории возмущений}
\exercise
\begin{center}
\textit{Найти смещение уровня энергии основного состояния атома водорода под влиянием конечных размеров ядра. Ядро считать равномерно заряженным по объему шаром радиуса $r_0$}
\end{center}

Для начала разберёмся, как повлияет конечный размер ядра на гамильтониан. Обычно мы считаем, что ядро является точечным зарядом, то есть $U(r) = -e^2/r$ на всем пространстве. В случае же равномерно заряженного по объему шара известного радиуса, потенциал будет следующим:
\[
U(r)=\begin{cases}
    -\frac{e^2}{r}, &r\in [r_0, +\infty) \\
    \frac{e^2}{2r_0}\left(\frac{r^2}{r_0^2}-3\right), &r\in(0, r_0)
\end{cases}
\]
Кратко обсудим то, откуда взялся этот потенциал. В общем случае потенциал в точке $\vec{r}$ с известным распределением заряда задается формулой $\varphi(\vec{r}) = \int_{\mathbb{R}^3}\rho(\vec{r})/|\vec{r}-\vec{r}'|d^3r$, где $\rho(\vec{r}')$ -- плотность заряда, распределенного внутри заряженного шара, $\vec{r}$ -- радиус-вектор точки наблюдения, $\vec{r}'$ -- радиус-вектор точки внутри заряженного шара. Тогда, считая заряд распределенным равномерно по всему объему ядра, мы можем найти плотность: $\rho(r)=3Q/4\pi r_0^3$. Далее, подставляя это значение в интеграл и переходя к сферическим координатам (так как в рамках данной задачи явно прослеживается сферическая симметрия), мы и получаем нашу итоговую формулу. Тогда, посчитаем наше возмущение:
\[
V(r) = \begin{cases}
    0, & r\in[r_0, +\infty) \\
    \frac{e^2}{2r_0}\left(\frac{r^2}{r_0^2}-3+\frac{2r_0}{r}\right), & r\in(0, r_0)
\end{cases}
\]
Видим, что оно отличается от нуля только внутри самого ядра. При этом, мы знаем характерный размер ядра водорода $r_0\sim 10^{-15}$м, а средний радиус орбиты электрона (боровский радиус) -- $a\sim 10^{-11}$м. То есть, возмущение, вызванное конечными размерами ядра водорода, действует на очень далеком расстоянии от электрона. Поэтому мы можем считать это возмущение малым и применить к нему теорию возмущений.

Для того чтобы посчитать смещение уровня энергии основного состояния атома водорода, нам нужно выписать волновую функцию этого состояния (мы подробно обсуждали ее вывод в $9$ семинаре прошлого семестра):
\[
\psi_{100}(r, \theta, \varphi) = \sqrt{\frac{1}{\pi a^3}}e^{-\frac{r}{a}},
\]
где $a$ -- радиус Бора.

Наконец, можем вычислить поправку по энергии, считая $r_0 \ll a$, то есть $e^{-\frac{r}{a}}\approx 1$, где $r\in(0, r_0)$:
\begin{align*}
    E_{100}^{(1)} &= \bra{\psi_{100}^{(0)}}\hat{V}\ket{\psi_{100}^{(0)}} = \frac{1}{\pi a^3} \frac{e^2}{2r_0} \int\limits_{0}^{2\pi}d\varphi\int\limits_{0}^{\pi}\sin\theta d\theta\int\limits_{0}^{r_0}r^2e^{-\frac{2r}{a}}\left(\frac{r^2}{r_0}-3+\frac{2r_0}{r}\right)dr = \\
    &= \frac{2e^2}{a^3r_0}\int\limits_0^{r_0}\left(\frac{r^4}{r_0}-3r^2+2rr_0\right)dr = \frac{2e^2}{a^3}\frac{r_0^2}{5} = -\frac{4}{5}\left(\frac{r_0}{a}\right)^2E_0^{(0)}
\end{align*}
Получается, поправка составляет примерно $10^{-8}$ часть от исходной энергии, что еще раз подтверждает, что теория возмущений в данной ситуации применима.
\csquare{darklavender}

\exercise
\begin{center}
\textit{Найти вероятность перехода между состояниями дискретного спектра под действием возмущений: 
$a) \; V(t) = V_0 e^{-t^2/\tau^2}$, $\textit{б}) \; V(t) = \frac{V_0}{2}\left(1+\frac{2}{\pi}\arctan\frac{t}{\tau}\right)$}
\end{center}

Так как мы работаем с дискретным спектром, для поиска вероятности перехода можем воспользоваться формулами из предыдущего семинара. Будем учитывать только поправки первого порядка, так как они вносят наибольший вклад. Тогда, вероятность выражается следующим образом:
\[
P(i\to n) = \left|c_n^{(1)}(t)\right|^2 = \left|-\frac{i}{\hbar}\int\limits_{t_0}^t e^{i\omega_{ni}t_1}\hat{V}(t_1)_{ni}dt_1\right|^2, \quad V_{ni} = \bra{n}\hat{V}\ket{i}, \; \omega_{ni} = \frac{E_n-E_i}{\hbar}.
\]

Так как нас интересует вероятность того, что переход произошел в принципе, без привязки к конкретному моменту времени, интегрировать будем в пределах от $-\infty$ до $+\infty$. 

$a) V(t) =  V_0 e^{-t^2/\tau^2}$. Подставим данное возмущение в формулу для нахождения вероятности:
\begin{align*}
    P(i\to n) = \frac{1}{\hbar^2}\left|\int\limits_{-\infty}^{+\infty}\bra{n}\hat{V}_0 e^{-t^2/\tau^2}\ket{i}e^{i\omega_{ni}t}dt\right|^2 = \frac{\left|V_{0ni}\right|^2}{\hbar^2}\left|\int\limits_{-\infty}^{+\infty}e^{-t^2/\tau^2+i\omega_{ni}t}dt\right|^2,
\end{align*}
где $V_{0ni} = \bra{n}\hat{V}_0\ket{i}, \; \omega_{ni} = (E_n-E_i)/\hbar$.

Последний интеграл путем замены переменной можно свести к интегралу Пуассона:
\[
\int\limits_{-\infty}^{+\infty}e^{-t^2/\tau^2+i\omega_{ni}t}dt = \sqrt{\pi}|\tau|e^{-\omega_{ni}^2\tau^2/4} \implies \left|\int\limits_{-\infty}^{+\infty}e^{-t^2/\tau^2+i\omega_{ni}t}dt\right|^2 = \pi\tau^2 e^{-\omega_{ni}^2\tau^2/2}.
\]

Итого, искомая вероятность
\[
P(n\to i) = \frac{\pi}{\hbar^2}\tau^2e^{-\omega_{ni}^2\tau^2/2}\left|V_{0ni}\right|^2
\]

$\textit{б}) \; V(t) = \frac{V_0}{2}\left(1+\frac{2}{\pi}\arctan\frac{t}{\tau}\right)$. Поступим аналогично предыдущему пункту:
\[
P(i\to n) = \frac{|V_{0ni}|^2}{4\hbar^2}\left|\int\limits_{-\infty}^{+\infty}\left(1+\frac{2}{\pi}\arctan\frac{t}{\tau}\right)e^{i\omega_{ni}t}dt\right|^2
\]

Рассмотрим этот интеграл подробнее. Видно, что напрямую его взять не получится, так как он расходится. Для того чтобы обойти эту проблему, проинтегрируем по частям:
\[
\int\limits_{-\infty}^{+\infty}\left(1+\frac{2}{\pi}\arctan\frac{t}{\tau}\right)e^{i\omega_{ni}t}dt = \left.\frac{e^{i\omega_{ni}t}}{i\omega_{ni}}\left(1+\frac{2}{\pi}\arctan\frac{t}{\tau}\right)\right|_{-\infty}^{+\infty}-\frac{2\tau}{i\pi\omega_{ni}}\int\limits_{-\infty}^{+\infty}\frac{e^{i\omega_{ni}t}}{t^2+\tau^2}dt.
\]

Первое слагаемое не влияет на вероятность, так как при вычислении квадрата модуля оно дает не более чем ограниченный вклад, и описывает поправку к волновой функции. Получается, нас интересует только второе слагаемое, интеграл в котором ищется через теорему о вычетах ($t=\pm i\tau$ -- простые полюса):
\[
\int\limits_{-\infty}^{+\infty}\frac{e^{i\omega_{ni}t}}{t^2+\tau^2}dt = \frac{\pi}{\tau}e^{-\omega_{ni}\tau}.
\]

Подставляя этот интеграл в формулу вероятности, найдем
\[
P(i\to n) = \frac{|V_{0ni}|^2}{\hbar^2\omega_{ni}^2}e^{-2\omega_{ni}\tau}.
\]
\csquare{darklavender}

\exercise
\begin{center}
\textit{Покоящийся мюон помещен в магнитное поле $\vec{\mathcal{H}}(t)$, медленно
прецессирующее вокруг оси $z$, составляя с ней постоянный угол $\theta$. Найдите зависимость от времени спиновой функции $\ket{\chi(t)}$ и поляризации $\vec{P}(t)$ мюона, если в начальный момент времени его состояние $\ket{\chi(0)}$ определяется положительной проекцией спина на направление магнитного поля. Покажите, что за один полный оборот вектора магнитного поля спиновая функция приобретает фазу, пропорциональную телесному углу, который охватывает вектор $\vec{\mathcal{H}}(t)$ (фаза Берри).}
\end{center}

Сначала нам надо понять, как выглядит наше магнитное поле. Мы знаем, что модуль $\vec{\mathcal{H}}$ не меняется, как и угол с осью $z$. Значит, в сферических координатах меняется только $\varphi(t)=\omega t$:
\[
\vec{\mathcal{H}}(t) = \mathcal{H}(\sin\theta\cos\omega t, \; \sin\theta\sin\omega t, \; \cos\theta)^T.
\]
В стандартном случае, когда $\vec{\mathcal{H}}||z$, решение выглядит так:
\[
\ket{+}=\begin{pmatrix} 1 \\ 0 \end{pmatrix}, \quad \ket{-}=\begin{pmatrix} 0 \\ 1 \end{pmatrix},
\]

а направляющий вектор поля -- так: $\vec{n}=(\sin\theta\cos\omega t, \; \sin\theta\sin\omega t, \; \cos\theta)^T$.

Посмотрим, как выглядит решение невозмущенной задачи, считая вектор магнитного поля фиксированным: $\vec{\mathcal{H}}(t) = \mathcal{H}(\sin\theta\cos\varphi, \; \sin\theta\sin\varphi, \; \cos\theta)^T.$ В прошлом семестре мы искали собственные значения и функции спинового оператора $\hat{\sigma}_n=(\hat{\vec{\sigma}},\vec{n})$. Решение этой задачи можно найти в 8 семинаре, а сейчас выпишем только готовые собственные векторы:
\begin{align*}
    \ket{+}_{\mathcal{H}} &= \cos\frac{\theta}{2}\ket{+}+\sin\frac{\theta}{2}e^{i\varphi}\ket{-}, \\
    \ket{-}_{\mathcal{H}} &= -\sin\frac{\theta}{2}\ket{+}+\cos\frac{\theta}{2}e^{i\varphi}\ket{-}.
    \end{align*}
Значения энергии, соответствующие состояниям $\ket{\pm}_{\mathcal{H}}$ от наличия прецессии не меняются: $E_{\pm}=\mp\mu_0\mathcal{H}/2$. 

Обсудим, что значит ``медленная прецессия'' магнитного поля. Мы можем считать, что гамильтониан зависит от некоторого параметра $\xi(t)$: $\hat{H}(t)=\hat{H}(\xi(t))$. Вследствие плавного, медленного изменения гамильтониана, вероятность квантового перехода системы между мгновенными собственными состояниями является экспоненциально малой, что значит, что система, в начальный момент находившаяся в состоянии, соответствующем невырожденному уровню энергии $E_n(\xi(0))$, будет и дальше находиться в состоянии, отвечающем тому же уровню энергии $E_n(\xi(t))$. Отсюда мы можем понять вид волновой функции: $\ket{\psi(t)}=e^{i\gamma(t)}\hat{U}(t)\ket{\psi(0)}$, где $\hat{U}(t) = e^{-iE_n t/\hbar}$ -- оператор эволюции, из-за действия которого волновая функция приобретает \textit{динамическую} фазу, а $e^{i\gamma(t)}$ -- \textit{геометрическая} фаза, появляющаяся вследствие геометрических свойств пространства параметров гамильтониана. Эта геометрическая фаза носит название \textit{фазы Берри}. В общем (одномерном) случае она считается по следующей формуле:
\[
\gamma(t)=i\int\limits_0^t dt'\bra{\psi(\xi(t'))}\frac{d}{dt'}\ket{\psi(\xi(t'))} = i\int\limits_{\xi(0)}^{\xi(t)}d\xi\bra{\psi(\xi)}\frac{d}{d\xi}\ket{\psi(\xi)}.
\]

Так как по условию нам нужно найти фазу, набежавшую за один полный оборот вектора магнитного поля, наш изменяющийся параметр $\xi = \varphi$ и меняется в промежутке от $0$ до $2\pi$:
\begin{align*}
    \gamma&=i\int\limits_0^{2\pi}d\varphi_{\mathcal{H}}\bra{+}\frac{\partial}{\partial\varphi}\ket{+}_{\mathcal{H}}=i\int\limits_0^{2\pi}\begin{pmatrix} \cos\frac{\theta}{2}&\sin\frac{\theta}{2}e^{-i\varphi} \end{pmatrix}\begin{pmatrix} 0\\i\sin\frac{\theta}{2}e^{i\varphi} \end{pmatrix}d\varphi = -\int\limits_0^{2\pi}\sin^2\frac{\theta}{2}d\varphi = \\
    &=-2\pi\sin^2\frac{\theta}{2} = -\pi(1-\cos\theta)
\end{align*}

Телесный угол, захватываемый вектором магнитного поля за один оборот равен $\Omega = 2\pi(1-\cos\theta)$, то есть фаза Берри равна его половине, взятой с минусом: $\gamma=-\Omega/2$.

Если рассматривать не один период, а произвольный промежуток времени, то фаза Берри равна
\[
\gamma(t)=i\int\limits_0^{\varphi(t)}d\varphi_{\mathcal{H}}\bra{+}\frac{\partial}{\partial\varphi}\ket{+}_{\mathcal{H}}=-\int\limits_0^{\varphi(t)}\sin^2\frac{\theta}{2}d\varphi = -\varphi(t)\sin^2\frac{\theta}{2}.
\]

Наконец, найдем спиновую функцию:
\[
\chi(t)=e^{-i\varphi(t)\sin^2\frac{\theta}{2}}e^{i\frac{\mu_0\mathcal{H}}{2\hbar}}\left(\cos\frac{\theta}{2}\ket{+}+\sin\frac{\theta}{2}e^{i\varphi(0)}\ket{-}\right).
\]
\csquare{darklavender}